\chapter{E-cigarettes}

\section*{Definition and Overview}
E-cigarettes, also called electronic cigarettes or vapes, are battery-powered devices that heat a liquid containing nicotine and other chemicals into an aerosol that is inhaled. They are promoted as a way to reduce the harm of smoking and as a quitting aid, but they are not harmless, and the use of e-cigarettes among young people has risen sharply. In Australia, nicotine vaping products are regulated and are only legally available with a prescription for people trying to quit smoking. This chapter explains how e-cigarettes work, their health effects, and what they mean for smokers and young people.

\section*{Epidemiology}
Vaping has increased dramatically in Australia, particularly among young people, with a substantial proportion of teenagers having tried e-cigarettes. While adult smoking rates have fallen, vaping among youth is a growing public health concern because of nicotine addiction and the health effects of aerosol inhalation. Among adults, e-cigarettes are used mainly by current or former smokers, and a minority of people who vape have never smoked. The debate continues over the balance of benefits for smokers and risks for non-smokers and young people.

\section*{Aetiology and Pathophysiology}
The effects of vaping depend on the contents of the liquid and the aerosol.
\begin{itemize}
    \item \textbf{Nicotine:} the addictive chemical in most vape liquids; it raises heart rate and blood pressure, and it is particularly harmful to the developing adolescent brain
    \item \textbf{Other chemicals:} aerosols contain propylene glycol, vegetable glycerin, flavourings, and heavy metals, which irritate the airways and may cause lung damage
    \item \textbf{Mechanism of harm:} inhaled particles and chemicals trigger airway inflammation, impair immune defences, and can cause coughing, wheezing, and lung injury
    \item \textbf{Addiction:} nicotine produces rapid dependence, and withdrawal causes irritability, cravings, and difficulty concentrating
    \item \textbf{Dual use:} many vapers also smoke, and the two habits may reinforce each other
    \item \textbf{E-cigarette or vaping-associated lung injury (EVALI):} a rare but serious acute lung illness linked to certain vaping products
\end{itemize}

\section*{Clinical Features}
The symptoms of vaping-related harm can be acute or chronic.
\begin{itemize}
    \item Cough, throat irritation, and a dry mouth
    \item Wheezing, chest tightness, and breathlessness
    \item Nausea, headaches, and dizziness, particularly with high nicotine intake
    \item Palpitations and raised heart rate
    \item Mouth and gum problems
    \item Symptoms of nicotine dependence: cravings, irritability, and difficulty stopping
    \item In EVALI: severe breathlessness, chest pain, cough, fever, and gastrointestinal symptoms
\end{itemize}

\section*{Differential Diagnosis}
Respiratory symptoms in people who vape may have other causes.
\begin{itemize}
    \item Asthma and other chronic lung conditions
    \item Respiratory infections, including bronchitis and pneumonia
    \item Smoking-related lung disease in dual users
    \item Allergy and hay fever
    \item Anxiety, which can cause breathlessness and palpitations
\end{itemize}

\section*{When to Seek Medical Care}
Anyone experiencing persistent cough, breathlessness, chest pain, or other symptoms related to vaping should see a doctor. People who want to quit smoking and are considering vaping should discuss the options and the safest approach with their GP.

\textbf{Call 000 (or local emergency services) for:}
\begin{itemize}
    \item Sudden severe breathlessness or chest pain
    \item Difficulty breathing or wheezing that does not settle
    \item Coughing up blood
    \item Collapse or severe dizziness
\end{itemize}

\section*{Diagnosis and Investigations}
Assessment focuses on the extent of use and any health effects.
\begin{itemize}
    \item \textbf{History:} what products are used, how often, nicotine content, and any symptoms
    \item \textbf{Examination:} blood pressure, heart rate, and lung examination
    \item \textbf{Lung function testing:} spirometry if respiratory symptoms are present
    \item \textbf{Chest imaging:} X-ray or CT if lung injury is suspected
    \item \textbf{Assessment of dependence:} screening for nicotine dependence and readiness to quit
\end{itemize}

\section*{Management}
Management depends on whether the person is a smoker trying to quit, a non-smoker using vapes, or a young person taking them up.
\begin{itemize}
    \item \textbf{Smoking cessation:} for smokers, e-cigarettes are one option among several, including nicotine replacement therapy and counselling; the goal is complete cessation of all nicotine use
    \item \textbf{Stopping vaping:} for non-smokers and young people, the aim is to stop; support includes counselling, gradual reduction, and nicotine replacement where needed
    \item \textbf{Treatment of symptoms:} management of respiratory symptoms and any lung injury
    \item \textbf{Regulation and advice:} in Australia, prescription-only access to nicotine vaping products, with pharmacies supplying them for smoking cessation
    \item \textbf{Support services:} Quitline and other services provide free, non-judgemental support
\end{itemize}

\section*{Complications}
\begin{itemize}
    \item Nicotine addiction, particularly in young people, with effects on the developing brain
    \item Respiratory symptoms and lung injury, including EVALI
    \item Increased risk of starting to smoke among young people who vape
    \item Cardiovascular effects from nicotine
    \item Injuries from device explosions, including burns
    \item Poisoning in young children from swallowing vape liquid
\end{itemize}

\section*{Prognosis and Outcomes}
The health effects of vaping are still emerging, but evidence shows that e-cigarettes are not harmless, particularly for young people and non-smokers. For adult smokers, switching completely to e-cigarettes is likely to be less harmful than continuing to smoke, and some smokers quit successfully with their help, but the safest option is to stop all nicotine use. Respiratory symptoms usually improve when vaping stops.

\section*{Prevention and Self-Care}
\begin{itemize}
    \item Never start vaping; it is not a safe alternative to not using nicotine at all
    \item If you smoke, discuss the full range of quitting options with your GP, including proven nicotine replacement therapies
    \item If you vape and want to stop, set a date, reduce gradually, and seek support from Quitline
    \item Store vape liquids out of reach of children
    \item Avoid sharing devices, which spreads infection
    \item Talk openly with young people about the risks of vaping
    \item Seek medical advice for any persistent respiratory symptoms
\end{itemize}

\section*{Sources and Further Reading}
The following reputable sources provide further information for healthcare students and patients seeking reliable, current advice about e-cigarettes.
\begin{itemize}
    \item Quitline and Quit Victoria. \textit{Vaping information and quitting support.} https://www.quit.org.au/
    \item Australian Government Department of Health and Aged Care. \textit{E-cigarettes and vaping information.} https://www.health.gov.au/
    \item Therapeutic Goods Administration. \textit{Regulation of nicotine vaping products.} https://www.tga.gov.au/
    \item NHS. \textit{E-cigarettes and vaping: health information.} https://www.nhs.uk/better-health/quit-smoking/vaping/
    \item Centers for Disease Control and Prevention. \textit{Electronic cigarettes and lung injury.} https://www.cdc.gov/
    \item World Health Organization. \textit{Tobacco and e-cigarettes fact sheets.} https://www.who.int/
\end{itemize}
